%%%%%%%%%%%%%%%%%%%%%%%%%%%%%%%%%%%%%%%%%%%%%%%%%%%%%%%%%%%%%%%%%%%%%%
% Java OOP Practical Cheatsheet — CS2030S PE format
% Matches cs2030s-midterms-cheatsheet-ivan.tex style
\documentclass[10pt]{article}
\usepackage[scaled=0.92]{helvet}
\usepackage{calc}
\usepackage{multicol}
\usepackage{ifthen}
\usepackage[a4paper,margin=6mm,landscape]{geometry}
\usepackage{amsmath,amsthm,amsfonts,amssymb}
\usepackage{color,graphicx,overpic}
\usepackage{hyperref}
\usepackage{newtxtext}
\usepackage{enumitem}
\usepackage{array}
\usepackage[table]{xcolor}
\usepackage{minted}
\usemintedstyle{tango}
\definecolor{newblue}{HTML}{ADD8E6}
\colorlet{lightblue}{newblue!45}
\pagestyle{empty}
\setlist{nosep}

\makeatletter
\renewcommand{\section}{\@startsection{section}{1}{0mm}%
                                {-0.5ex plus -.3ex minus -.2ex}%
                                {0.3ex plus .2ex}%
                                {\normalfont\large\bfseries}}
\renewcommand{\subsection}{\@startsection{subsection}{2}{0mm}%
                                {-0.5ex plus -.3ex minus -.2ex}%
                                {0.3ex plus .2ex}%
                                {\normalfont\normalsize\bfseries}}
\renewcommand{\subsubsection}{\@startsection{subsubsection}{3}{0mm}%
                                {-0.5ex plus -.3ex minus -.2ex}%
                                {0.5ex plus .2ex}%
                                {\normalfont\small\bfseries}}
\renewcommand{\familydefault}{\sfdefault}
\renewcommand\rmdefault{\sfdefault}
\renewcommand{\labelenumii}{\theenumii}
\renewcommand{\theenumii}{\theenumi.\arabic{enumii}.}
\renewcommand\labelitemii{•}
\makeatother

\definecolor{myblue}{cmyk}{1,.72,0,.38}
\everymath\expandafter{\the\everymath \color{myblue}}
\setcounter{secnumdepth}{0}
\setlength{\parindent}{0pt}
\setlength{\parskip}{0pt plus 0.3ex}
\setlength{\leftmargini}{0.5cm}
\setlength{\leftmarginii}{0.5cm}
\setlist[itemize,1]{leftmargin=1.5mm,labelindent=0.5mm,labelsep=0.5mm}
\setlist[itemize,2]{leftmargin=3mm,labelindent=0.5mm,labelsep=0.5mm}

\begin{document}
\raggedright
\footnotesize
\begin{multicols}{3}

\setlength{\columnseprule}{0.25pt}
\setlength{\premulticols}{1pt}
\setlength{\postmulticols}{1pt}
\setlength{\multicolsep}{1pt}
\setlength{\columnsep}{1.5pt}

\begin{center}
    \fbox{%
        \parbox{0.6\linewidth}{\centering \textcolor{black}{
            {\normalsize\textbf{CS2030S PE}}
            \\ \small{Java OOP Practical}}
            \\ {\scriptsize \textcolor{myblue}{by @lostmusician}}
        }%
    }
\end{center}

\section{Terminal \& Testing}
\begin{minted}[breaklines,breakanywhere,bgcolor=lightblue]{bash}
# Run with input file
java Main < inputs/PE1.1.in

# Redirect output to file, then vim-diff
java PE1 < inputs/PE1.1.in > OUT
vim -d OUT outputs/PE1.1.out

# Inline diff (shows differences directly)
diff <(java Main < inputs/PE1.1.in) outputs/PE1.1.out
\end{minted}

\section{Common Signatures}
\begin{minted}[breaklines,breakanywhere,bgcolor=lightblue]{java}
public static void main(String[] args)
public boolean equals(Object o)
public int compareTo(Type o)  // Comparable
public int hashCode()
\end{minted}

\section{Equality \& Comparison}
\begin{minted}[breaklines,breakanywhere,bgcolor=lightblue]{java}
@Override
public boolean equals(Object o) {
  if (this == o) return true;
  if (o == null || getClass() != o.getClass()) return false;
  ClassName that = (ClassName) o;
  return n == that.n &&
    (s == null ? that.s == null : s.equals(that.s));
}

// Comparable<ClassName>
@Override
public int compareTo(ClassName o) {
  if (n != o.n) return n < o.n ? -1 : 1;
  if (s == null && o.s == null) return 0;
  if (s == null) return -1;
  if (o.s == null) return 1;
  return s.compareTo(o.s);
}
\end{minted}

\section{Event Handling}
\subsection{Event Pattern}
Each event extends \texttt{Event}, stores context, and returns the \textbf{next} events from \texttt{simulate()}.
\begin{minted}[breaklines,breakanywhere,bgcolor=lightblue]{java}
class TicketSubmittedEvent extends Event {
  private RTSystem rtsys;
  private Ticket ticket;

  public TicketSubmittedEvent(double time,
      RTSystem rtsys, Ticket ticket) {
    super(time);          // pass time to Event
    this.rtsys = rtsys;
    this.ticket = ticket;
  }

  @Override
  public Event[] simulate() {
    try {
      this.rtsys.fileTicket(ticket);
    } catch (CongestionException e) {
      return new Event[] {
        new TicketDroppedEvent(this.getTime(), e, ticket)
      };
    }
    return new Event[] {};   // no follow-up events
  }

  @Override
  public String toString() {
    return super.toString() + " " + this.ticket
        + " submitted " + this.rtsys;
  }
}
\end{minted}

\subsection{Event Checklist}
\begin{itemize}
  \item Always \textbf{call \texttt{super(time)}} as first line of constructor.
  \item \texttt{simulate()} returns \texttt{Event[]} — next events to schedule; return \texttt{new Event[]\{\}} if none.
  \item Use \texttt{this.getTime()} (inherited) to pass time to follow-up events.
  \item Catch \textbf{checked exceptions} inside \texttt{simulate()} and return error events; let \textbf{unchecked} propagate.
  \item \texttt{toString()} should chain \texttt{super.toString()} for the timestamp prefix.
  \item Keep event classes small — delegate logic to domain objects (\texttt{rtsys.fileTicket(...)}).
\end{itemize}

\section{Queue \& Inheritance}
\subsection{Queue Initialisation}
\begin{minted}[breaklines,breakanywhere,bgcolor=lightblue]{java}
// Basic queue of fixed size
Queue<Integer> q = new Queue<>(10);

// Array of generic queues (unchecked cast required)
@SuppressWarnings("unchecked")
Queue<Passenger>[] stops =
    (Queue<Passenger>[]) new Queue<?>[totalStops];
for (int i = 0; i < totalStops; i++)
    stops[i] = new Queue<>(capacity);

// Queue of queues
Queue<Queue<Task>> qOfQ = new Queue<>(5);
\end{minted}

\subsection{Queue Implementation}
\begin{minted}[breaklines,breakanywhere,bgcolor=lightblue]{java}
class Queue<T> {
  private T[] items;
  private int first;
  private int last;
  private int maxSize;
  private int len;

  public Queue(int size) {
    this.maxSize = size;
    this.first = -1;
    this.last = -1;
    this.len = 0;
    @SuppressWarnings("unchecked")
    T[] temp = (T[]) new Object[size];
    this.items = temp;
  }

  public boolean enq(T e) {
    if (this.isFull()) return false;
    if (this.isEmpty()) { this.first = 0; this.last = 0; }
    else { this.last = (this.last + 1) % this.maxSize; }
    this.items[last] = e;
    this.len += 1;
    return true;
  }

  public T deq() {
    if (this.isEmpty()) return null;
    T item = this.items[this.first];
    this.first = (this.first + 1) % this.maxSize;
    this.len -= 1;
    return item;
  }

  public T peek() {
    if (this.isEmpty()) return null;
    return this.items[this.first];
  }

  public boolean isFull()  { return this.len == this.maxSize; }
  public boolean isEmpty() { return this.len == 0; }
  public int length()      { return this.len; }

  @Override
  public String toString() {
    String str = "[ ";
    int i = this.first;
    int count = 0;
    while (count < this.len) {
      str += this.items[i] + " ";
      i = (i + 1) % this.maxSize;
      count++;
    }
    return str + "]";
  }
}
\end{minted}

\subsection{enq Circular Logic}
\begin{enumerate}
  \item If \texttt{isFull()}, return \texttt{false}.
  \item If \texttt{isEmpty()}, set \texttt{first = last = 0}.
  \item Otherwise, \texttt{last = (last + 1) \% maxSize}.
  \item Store element at \texttt{items[last]}, increment \texttt{len}, return \texttt{true}.
\end{enumerate}

\subsection{Comparable Rules}
\begin{itemize}
  \item \texttt{@Override public int compareTo(T other)}
  \item Returns \texttt{< 0} if \texttt{this} smaller; \texttt{0} equal; \texttt{> 0} larger
  \item Compare queues by \texttt{length()}; strings via \texttt{s.compareTo(t)}
\end{itemize}

\section{Generics \& Wildcards (PECS)}
\subsection{Generic Class \& Method}
\begin{minted}[breaklines,breakanywhere,bgcolor=lightblue]{java}
class Box<T> {
  private T v;
  public Box(T v) { this.v = v; }
  public T get()  { return v; }
  public <U> Box<U> map(Transformer<T,U> f) {
    return new Box<>(f.transform(v));
  }
}

// Bounded: T must be Number or subclass
class NumBox<T extends Number> {
  private T value;
  public double getDouble() { return value.doubleValue(); }
}
// Multiple bounds: <T extends Shape & Drawable>
\end{minted}

\subsection{PECS}
\textbf{P}roducer \textbf{E}xtends, \textbf{C}onsumer \textbf{S}uper
\begin{minted}[breaklines,breakanywhere,bgcolor=lightblue]{java}
// ? extends T: READING (producer — getters)
void read(Seq<? extends Shape> src) {
  Shape s = src.get(0); // OK
  // src.add(new Circle()); // CANNOT add
}

// ? super T: WRITING (consumer — setters)
void write(Seq<? super Circle> dest) {
  dest.add(new Circle(5.0)); // OK
  // Circle c = dest.get(0); // CANNOT get specific type
}
\end{minted}
\begin{itemize}
  \item \texttt{<?>}: Unbounded (unknown type)
  \item \texttt{<? extends T>}: Upper bounded; \textbf{read-only}
  \item \texttt{<? super T>}: Lower bounded; \textbf{write-safe}
\end{itemize}

\subsection{Generic Arrays}
Cannot use \texttt{new T[]} (type erasure). Cast from \texttt{Object[]}:
\begin{minted}[breaklines,breakanywhere,bgcolor=lightblue]{java}
@SuppressWarnings("unchecked")
Queue<Passenger>[] temp = (Queue<Passenger>[]) new Queue<?>[totalStops];
\end{minted}
Suppress only after justifying: array is \texttt{private}; only modified by typed setter; only returns subtypes of T.

\subsection{Type Erasure}
After compilation, generic types replaced by upper bound (\texttt{Object}). Cannot use \texttt{new T()}, \texttt{new T[]}, \texttt{T.class}.

\section{Exceptions \& Factory Methods}
\subsection{Exception Hierarchy}
\begin{itemize}
  \item \textbf{Unchecked}: \texttt{RuntimeException} — programmer error; no \texttt{throws} needed
  \item \textbf{Checked}: \texttt{Exception} — must \texttt{throws} or \texttt{try-catch}
\end{itemize}
\begin{minted}[breaklines,breakanywhere,bgcolor=lightblue]{java}
try {
  int result = riskyOperation();
} catch (SpecificException e) {
  // handle specific
} catch (Exception e) {
  // handle general
} finally {
  // always runs (cleanup)
}
// Checked: declare in signature
void readFile() throws IOException { ... }
\end{minted}

\subsection{Custom Exceptions}
\begin{minted}[breaklines,breakanywhere,bgcolor=lightblue]{java}
// Checked (extends Exception)
public class InvalidOperandException extends Exception {
  public InvalidOperandException() {
    super("Invalid operand provided for operation");
  }
  public InvalidOperandException(String msg) { super(msg); }
}

// Unchecked (extends RuntimeException)
public class CongestionException extends RuntimeException {
  public CongestionException(String msg) { super(msg); }
}

// Throwing and catching
void process(int x) throws InvalidOperandException {
  if (x < 0) throw new InvalidOperandException();
}
try {
  process(-1);
} catch (InvalidOperandException e) {
  System.out.println(e.getMessage());
}
\end{minted}

\subsection{Seq + Exception Example}
\begin{minted}[breaklines,breakanywhere,bgcolor=lightblue]{java}
class RTSystem {
  private Seq<TicketQueue> seq;

  public RTSystem(int numQueues, int queueSize) {
    this.seq = new Seq<>(numQueues);
    for (int i = 0; i < numQueues; i++) {
      this.seq.set(i, new TicketQueue(i, queueSize));
    }
  }

  public RTSystem fileTicket(Ticket t)
      throws CongestionException {
    TicketQueue q = this.seq.get(t.getPriority());
    if (q.isFull()) throw new CongestionException(t.getPriority());
    q.enq(t);
    return this;
  }
}
\end{minted}

\section{Useful Patterns}
\subsection{Null-safe \& parse helpers}
\begin{minted}[breaklines,breakanywhere,bgcolor=lightblue]{java}
public String getSafe() { return (s == null) ? "" : s; }

int parseIntOr(String s, int d) {
  try { return Integer.parseInt(s); }
  catch (NumberFormatException e) { return d; }
}
\end{minted}

\section{Array Utilities}
\begin{minted}[breaklines,breakanywhere,bgcolor=lightblue]{java}
int sum(int[] a) {
  int s = 0;
  for (int i = 0; i < a.length; i++) s += a[i];
  return s;
}
double avg(int[] a) {
  return a.length == 0 ? 0.0 : (double) sum(a) / a.length;
}
int minIndex(int[] a) {
  int mi = 0;
  for (int i = 1; i < a.length; i++)
    if (a[i] < a[mi]) mi = i;
  return mi;
}
int[] copyOf(int[] a, int n) {
  int[] b = new int[n];
  int m = a.length < n ? a.length : n;
  for (int i = 0; i < m; i++) b[i] = a[i];
  return b;
}
// Insertion sort (stable, good for small arrays)
for (int i = 1; i < n; i++) {
  int v = a[i], j = i;
  while (j > 0 && a[j-1] > v) { a[j] = a[j-1]; j--; }
  a[j] = v;
}
\end{minted}

\section{OOP Design Quick Tips}
\begin{itemize}
  \item \textbf{SRP}: Each class has one role. Split god objects.
  \item \textbf{Encapsulate}: private fields; expose minimal API; return copies for mutable arrays.
  \item \textbf{Composition $>$ Inheritance}: unless clearly is-a.
  \item \textbf{Immutability}: \texttt{final} fields, no setters where possible.
  \item \textbf{equals/hashCode}: override both if used in sets/maps.
  \item \textbf{Tell, Don't Ask}: \texttt{circle.isLargerThan(5)} not \texttt{circle.getRadius() > 5}.
  \item \textbf{Avoid static overuse}: global state hurts testability.
  \item \textbf{Silent swallow}: never \texttt{catch(Exception e) \{\}}; log or rethrow.
\end{itemize}

\section{Access Modifiers}
\begin{itemize}
  \item \textbf{public}: everywhere; \textbf{protected}: package + subclasses; \textbf{(default)}: package only; \textbf{private}: within class
  \item \textbf{static}: class-level (shared); \textbf{final}: cannot reassign; \textbf{final class}: cannot extend; \textbf{final method}: cannot override
  \item Override: same signature, covariant return OK, cannot reduce visibility
\end{itemize}

\section{Debug Checklist}
\begin{itemize}
  \item Compile error? Fix first error, then re-run.
  \item Off-by-one? Check loop bounds and indices.
  \item \texttt{NullPointerException}: which var is null? Add null guards.
  \item \texttt{ArrayIndexOutOfBounds}: check \texttt{arr.length} and indices.
  \item Wrong output: trace with \texttt{System.out.println} intermediate values.
\end{itemize}

\section{String.format Guide}
\begin{itemize}
  \item \texttt{\%s} string, \texttt{\%d} int, \texttt{\%f} float, \texttt{\%b} bool, \texttt{\%c} char, \texttt{\%x} hex
  \item Width: \texttt{\%8d}; precision: \texttt{\%.2f}; both: \texttt{\%8.2f}
  \item Flags: \texttt{-} left align, \texttt{0} zero-pad, \texttt{+} show sign, \texttt{,} grouping
\end{itemize}
\begin{minted}[breaklines,breakanywhere,bgcolor=lightblue]{java}
String.format("%s scored %d/%d (%.1f%%)",
    name, got, total, got*100.0/total)
String.format("Hex: 0x%08X", n)
String.format("%-10s | %8.2f", item, price)
\end{minted}

\end{multicols}
\end{document}
